% Generated. Do not edit.
% Projection scriptural-date-location/section, version 1: the Scriptural Date and
% Location page, generated whole from research/chronology.toml, the Scripture
% chronology corpus and the source library by tools/tpt document-types refresh.
% document: liturgy/roman-rite/postconciliar/roman-missal-third-edition-en-us-2011/propers/temporal/pc-s51-twenty-fifth-sunday-in-ordinary-time-year-a
\runninghead{Scriptural Date and Location}
\section{Scriptural Date and Location}
\label{sec:date-location}

{\small Each passage the Mass reads or sings from Scripture, in canonical order. The Context column names what is dated --- the event the passage narrates, a setting it remembers, the traditional attribution of its author, or the composition of its book --- and where it belongs; the Date column gives the answer of the Scripture chronology corpus for that subject, with every alternative it records. Texts of the Mass that are not drawn from Scripture have no row. The evidence for each date follows the table.\par}
\medskip
% Generated. Do not edit.
% Deterministic projection of research/chronology.toml.
% schema: 5
% document: liturgy/roman-rite/postconciliar/roman-missal-third-edition-en-us-2011/propers/temporal/pc-s51-twenty-fifth-sunday-in-ordinary-time-year-a
% generated-by: tools/tpt proper-chronology annotations
\newcommand{\chronologyannotationclaim}[6]{#6}
\newcommand{\chronologyannotationcomparisonclaim}[8]{#8}
\newcommand{\chronologyannotationreach}[3]{}
\newcommand{\chronologyannotationgroup}[3]{#3}
\newcommand{\chronologyannotationcomparisongroup}[4]{#4}
\newcommand{\chronologyannotation}[1]{%
  \ifcsname triptychchronologyannotation@#1\endcsname
    \csname triptychchronologyannotation@#1\endcsname
  \else
    \PackageError{triptych}{No chronology annotation for #1}{}%
  \fi}
% appointment-dependency: guidance/liturgy/postconciliar-propers-registry.md sha256=376860a7c37a4c75bb6d880cd1b260183be1877a07e3cd01660a9e2854e2ea72
% appointment-dependency: src/gpt/liturgy/roman-rite/postconciliar/roman-missal-third-edition-en-us-2011/propers/registry/README.md sha256=df0aae579b2ae7e02cca5597f382ee615974bc84fc46fea909a9c21d942ce152
% appointment-dependency: src/gpt/liturgy/roman-rite/postconciliar/roman-missal-third-edition-en-us-2011/propers/registry/formula-dispositions.md sha256=af014f05723fb53a99006ef740101dc8abf629ed8b0b53a60069e94713cb289e
% appointment-dependency: src/gpt/liturgy/roman-rite/postconciliar/roman-missal-third-edition-en-us-2011/propers/temporal/pc-s51-twenty-fifth-sunday-in-ordinary-time-year-a/instance/manifest.md sha256=dfc0ceb23399512805805bf2c6f589ee80233218786414b9f888c985360dfb47
% appointment-dependency: src/gpt/liturgy/roman-rite/postconciliar/roman-missal-third-edition-en-us-2011/propers/temporal/pc-s51-twenty-fifth-sunday-in-ordinary-time-year-a/propers/verified.md sha256=e93ce380cc3b5f124cf1581b3854ed3321217552db8258791823832c9b2246a5
% appointment-dependency: src/gpt/liturgy/roman-rite/postconciliar/roman-missal-third-edition-en-us-2011/propers/temporal/pc-s51-twenty-fifth-sunday-in-ordinary-time-year-a/research/chronology-inputs.toml sha256=e0cc18992a9f9bd410546fad2bf9d8596535d32fde724f1c575906ab54616df1
% appointment-dependency: src/gpt/liturgy/roman-rite/postconciliar/roman-missal-third-edition-en-us-2011/propers/temporal/shared/ordinary-time/weeks/25/propers/verified.md sha256=ac29dbdf01322b68e4aa24547373f83636eb97a4b0f29de5837cb14a773aca8c
% appointment-note: entrance: Psalm 37:39-40; identified-basis; hebrew numbering; canonical verse queries, not appointed wording.
% appointment-note: collect: Composed Missal oration; no directly appointed biblical citation.
% appointment-note: first-reading: Isaiah 55:6-9; appointed; vulgate numbering; canonical verse queries, not appointed wording.
% appointment-note: responsorial-psalm: Psalm 145:2-3, 8-9, 17-18; appointed; hebrew numbering; canonical verse queries, not appointed wording.
% appointment-note: responsorial-psalm: Psalm 145:18a; response; hebrew numbering; whole-verse query envelope, not subverse chronology or appointed wording.
% appointment-note: second-reading: Philippians 1:20c-24, 27a; appointed; vulgate numbering; whole-verse query envelope, not subverse chronology or appointed wording.
% appointment-note: acclamation: Cf. Acts 16:14b; adaptation; vulgate numbering; whole-verse query envelope, not subverse chronology or appointed wording.
% appointment-note: gospel: Matthew 20:1-16a; appointed; vulgate numbering; whole-verse query envelope, not subverse chronology or appointed wording.
% appointment-note: offerings: Composed Missal oration; no directly appointed biblical citation.
% appointment-note: communion-ps119: Psalm 119:4-5; appointed; hebrew numbering; canonical verse queries, not appointed wording.
% appointment-note: communion-jn10: John 10:14; appointed; vulgate numbering; canonical verse queries, not appointed wording.
% appointment-note: after-communion: Composed Missal oration; no directly appointed biblical citation.
% comparison-dependency: src/gpt/liturgy/roman-rite/postconciliar/roman-missal-third-edition-en-us-2011/propers/temporal/pc-s51-twenty-fifth-sunday-in-ordinary-time-year-a/research/chronology-profile-comparisons.toml sha256=a39e84965db1fc5b391c2eb3d1923b61ae7b144b947eb9da9ab70031d3358c30
% entrance: Entrance antiphon
\expandafter\def\csname triptychchronologyannotation@entrance\endcsname{%
  \chronologyannotationgroup{historical-setting}{research-pending}{\textbf{Historical setting} -- Occasion date unresolved.}%
  \space \chronologyannotationgroup{composition}{preferred}{\textbf{Composition}: \chronologyannotationreach{Ps.36.39}{Ps}{true}\chronologyannotationreach{Ps.36.40}{Ps}{true}\chronologyannotationclaim{critical.psalms.latest-composition-boundary}{composition}{catholic-critical-v1}{preferred}{before the Maccabean period, around 165 B.C.; no individual psalm can be dated securely}{Before c. 165 B.C.}}%
}
% first-reading: First reading
\expandafter\def\csname triptychchronologyannotation@first-reading\endcsname{%
  \chronologyannotationgroup{traditional-attribution}{preferred}{\textbf{Traditional attribution}: \chronologyannotationreach{Is.55.6}{Is}{true}\chronologyannotationreach{Is.55.7}{Is}{true}\chronologyannotationreach{Is.55.8}{Is}{true}\chronologyannotationreach{Is.55.9}{Is}{true}\chronologyannotationclaim{israel.prophets.isaias-traditional-era}{traditional-attribution}{catholic-traditional-v1}{preferred}{from the closing year of Ozias, King of Juda (740 B.C.); no prophecies appear to be later than 701 B.C.}{Isaias (ministry in Souvay's traditional account), B.C. 740--701}.}%
  \space \chronologyannotationgroup{prophecy-given}{preferred}{\textbf{Prophecy given}: \chronologyannotationreach{Is.55.6}{Is.40 Is.41 Is.42 Is.43 Is.44 Is.45 Is.46 Is.47 Is.48 Is.49 Is.50 Is.51 Is.52 Is.53 Is.54 Is.55}{true}\chronologyannotationreach{Is.55.7}{Is.40 Is.41 Is.42 Is.43 Is.44 Is.45 Is.46 Is.47 Is.48 Is.49 Is.50 Is.51 Is.52 Is.53 Is.54 Is.55}{true}\chronologyannotationreach{Is.55.8}{Is.40 Is.41 Is.42 Is.43 Is.44 Is.45 Is.46 Is.47 Is.48 Is.49 Is.50 Is.51 Is.52 Is.53 Is.54 Is.55}{true}\chronologyannotationreach{Is.55.9}{Is.40 Is.41 Is.42 Is.43 Is.44 Is.45 Is.46 Is.47 Is.48 Is.49 Is.50 Is.51 Is.52 Is.53 Is.54 Is.55}{true}\chronologyannotationclaim{israel.exile.second-isaias-oracles}{prophecy-given}{catholic-critical-v1}{preferred}{toward the end of the Babylonian exile}{The prophetic activity behind Isaiah 40-55 in the NABRE introduction, toward the end of the Babylonian exile}.}%
}
% responsorial-psalm: Responsorial psalm
\expandafter\def\csname triptychchronologyannotation@responsorial-psalm\endcsname{%
  \chronologyannotationgroup{traditional-attribution}{disputed}{\textbf{Traditional attribution} -- disputed: \chronologyannotationreach{Ps.144.17}{Ps.107 Ps.144}{true}\chronologyannotationreach{Ps.144.18}{Ps.107 Ps.144}{true}\chronologyannotationreach{Ps.144.2}{Ps.107 Ps.144}{true}\chronologyannotationreach{Ps.144.3}{Ps.107 Ps.144}{true}\chronologyannotationreach{Ps.144.8}{Ps.107 Ps.144}{true}\chronologyannotationreach{Ps.144.9}{Ps.107 Ps.144}{true}\chronologyannotationclaim{israel.monarchy.david-traditional-era}{traditional-attribution}{catholic-traditional-v1}{disputed}{reigned from 1055 to 1015 B.C.}{David (reign in the usual chronology), B.C. 1055--1015}.}%
  \space \chronologyannotationgroup{historical-setting}{research-pending}{\textbf{Historical setting} -- Occasion date unresolved.}%
  \space \chronologyannotationgroup{composition}{preferred}{\textbf{Composition}: \chronologyannotationreach{Ps.144.17}{Ps}{true}\chronologyannotationreach{Ps.144.18}{Ps}{true}\chronologyannotationreach{Ps.144.2}{Ps}{true}\chronologyannotationreach{Ps.144.3}{Ps}{true}\chronologyannotationreach{Ps.144.8}{Ps}{true}\chronologyannotationreach{Ps.144.9}{Ps}{true}\chronologyannotationclaim{critical.psalms.latest-composition-boundary}{composition}{catholic-critical-v1}{preferred}{before the Maccabean period, around 165 B.C.; no individual psalm can be dated securely}{Before c. 165 B.C.}}%
}
% second-reading: Second reading
\expandafter\def\csname triptychchronologyannotation@second-reading\endcsname{%
  \chronologyannotationgroup{composition}{disputed}{\textbf{Composition} -- disputed: \chronologyannotationreach{Phil.1.20}{Phil}{true}\chronologyannotationreach{Phil.1.21}{Phil}{true}\chronologyannotationreach{Phil.1.22}{Phil}{true}\chronologyannotationreach{Phil.1.23}{Phil}{true}\chronologyannotationreach{Phil.1.24}{Phil}{true}\chronologyannotationreach{Phil.1.27}{Phil}{true}\chronologyannotationclaim{composition.epistle-to-the-philippians}{composition}{catholic-traditional-v1}{disputed}{(Philemon; Colossians; Ephesians; Philippians), 61}{A.D. 61}; \chronologyannotationreach{Phil.1.20}{Phil}{true}\chronologyannotationreach{Phil.1.21}{Phil}{true}\chronologyannotationreach{Phil.1.22}{Phil}{true}\chronologyannotationreach{Phil.1.23}{Phil}{true}\chronologyannotationreach{Phil.1.24}{Phil}{true}\chronologyannotationreach{Phil.1.27}{Phil}{true}\chronologyannotationclaim{composition.epistle-to-the-philippians}{composition}{catholic-traditional-v1}{disputed}{at Rome (A.D. 62-64)}{A.D. 62--64}.}%
}
% acclamation: Gospel acclamation
\expandafter\def\csname triptychchronologyannotation@acclamation\endcsname{%
  \chronologyannotationgroup{narrated-event}{disputed}{\textbf{Event} -- disputed: \chronologyannotationreach{Acts.16.14}{Acts.15.36- Acts.16 Acts.17 Acts.18.1-22}{true}\chronologyannotationclaim{apostolic-age.saint-paul-second-mission}{narrated-event}{catholic-traditional-v1}{disputed}{second mission, 50-53}{St Paul's Second Missionary Journey, A.D. 50--53}; \chronologyannotationreach{Acts.16.14}{Acts.15.36- Acts.16 Acts.17 Acts.18.1-22}{true}\chronologyannotationclaim{apostolic-age.saint-paul-second-mission}{narrated-event}{catholic-traditional-v1}{disputed}{began about three years previously, in 51 ... probably in 54}{St Paul's Second Missionary Journey, A.D. 51--54}.}%
  \space \chronologyannotationgroup{composition}{preferred}{\textbf{Composition}: \chronologyannotationreach{Acts.16.14}{Acts}{true}\chronologyannotationclaim{composition.acts-of-the-apostles}{composition}{catholic-traditional-v1}{preferred}{the year 64 is the most probable date}{c. A.D. 64}.}%
}
% gospel: Gospel
\expandafter\def\csname triptychchronologyannotation@gospel\endcsname{%
  \chronologyannotationgroup{narrated-event}{research-pending}{\textbf{Event} -- Narrated event date unresolved.}%
  \space \chronologyannotationgroup{composition}{disputed}{\textbf{Composition} -- disputed: \chronologyannotationreach{Matt.20.1}{Matt}{true}\chronologyannotationreach{Matt.20.10}{Matt}{true}\chronologyannotationreach{Matt.20.11}{Matt}{true}\chronologyannotationreach{Matt.20.12}{Matt}{true}\chronologyannotationreach{Matt.20.13}{Matt}{true}\chronologyannotationreach{Matt.20.14}{Matt}{true}\chronologyannotationreach{Matt.20.15}{Matt}{true}\chronologyannotationreach{Matt.20.16}{Matt}{true}\chronologyannotationreach{Matt.20.2}{Matt}{true}\chronologyannotationreach{Matt.20.3}{Matt}{true}\chronologyannotationreach{Matt.20.4}{Matt}{true}\chronologyannotationreach{Matt.20.5}{Matt}{true}\chronologyannotationreach{Matt.20.6}{Matt}{true}\chronologyannotationreach{Matt.20.7}{Matt}{true}\chronologyannotationreach{Matt.20.8}{Matt}{true}\chronologyannotationreach{Matt.20.9}{Matt}{true}\chronologyannotationclaim{composition.gospel-of-matthew}{composition}{catholic-traditional-v1}{disputed}{about A.D. 38-45}{c. A.D. 38--45}; \chronologyannotationreach{Matt.20.1}{Matt}{true}\chronologyannotationreach{Matt.20.10}{Matt}{true}\chronologyannotationreach{Matt.20.11}{Matt}{true}\chronologyannotationreach{Matt.20.12}{Matt}{true}\chronologyannotationreach{Matt.20.13}{Matt}{true}\chronologyannotationreach{Matt.20.14}{Matt}{true}\chronologyannotationreach{Matt.20.15}{Matt}{true}\chronologyannotationreach{Matt.20.16}{Matt}{true}\chronologyannotationreach{Matt.20.2}{Matt}{true}\chronologyannotationreach{Matt.20.3}{Matt}{true}\chronologyannotationreach{Matt.20.4}{Matt}{true}\chronologyannotationreach{Matt.20.5}{Matt}{true}\chronologyannotationreach{Matt.20.6}{Matt}{true}\chronologyannotationreach{Matt.20.7}{Matt}{true}\chronologyannotationreach{Matt.20.8}{Matt}{true}\chronologyannotationreach{Matt.20.9}{Matt}{true}\chronologyannotationclaim{composition.gospel-of-matthew}{composition}{catholic-traditional-v1}{disputed}{about the year 40-42}{c. A.D. 40--42}; \chronologyannotationreach{Matt.20.1}{Matt}{true}\chronologyannotationreach{Matt.20.10}{Matt}{true}\chronologyannotationreach{Matt.20.11}{Matt}{true}\chronologyannotationreach{Matt.20.12}{Matt}{true}\chronologyannotationreach{Matt.20.13}{Matt}{true}\chronologyannotationreach{Matt.20.14}{Matt}{true}\chronologyannotationreach{Matt.20.15}{Matt}{true}\chronologyannotationreach{Matt.20.16}{Matt}{true}\chronologyannotationreach{Matt.20.2}{Matt}{true}\chronologyannotationreach{Matt.20.3}{Matt}{true}\chronologyannotationreach{Matt.20.4}{Matt}{true}\chronologyannotationreach{Matt.20.5}{Matt}{true}\chronologyannotationreach{Matt.20.6}{Matt}{true}\chronologyannotationreach{Matt.20.7}{Matt}{true}\chronologyannotationreach{Matt.20.8}{Matt}{true}\chronologyannotationreach{Matt.20.9}{Matt}{true}\chronologyannotationclaim{composition.gospel-of-matthew}{composition}{catholic-traditional-v1}{disputed}{the years 40-45}{A.D. 40--45}; \chronologyannotationreach{Matt.20.1}{Matt}{true}\chronologyannotationreach{Matt.20.10}{Matt}{true}\chronologyannotationreach{Matt.20.11}{Matt}{true}\chronologyannotationreach{Matt.20.12}{Matt}{true}\chronologyannotationreach{Matt.20.13}{Matt}{true}\chronologyannotationreach{Matt.20.14}{Matt}{true}\chronologyannotationreach{Matt.20.15}{Matt}{true}\chronologyannotationreach{Matt.20.16}{Matt}{true}\chronologyannotationreach{Matt.20.2}{Matt}{true}\chronologyannotationreach{Matt.20.3}{Matt}{true}\chronologyannotationreach{Matt.20.4}{Matt}{true}\chronologyannotationreach{Matt.20.5}{Matt}{true}\chronologyannotationreach{Matt.20.6}{Matt}{true}\chronologyannotationreach{Matt.20.7}{Matt}{true}\chronologyannotationreach{Matt.20.8}{Matt}{true}\chronologyannotationreach{Matt.20.9}{Matt}{true}\chronologyannotationclaim{composition.gospel-of-matthew}{composition}{catholic-traditional-v1}{disputed}{about the year 60-68}{c. A.D. 60--68}; \chronologyannotationreach{Matt.20.1}{Matt}{true}\chronologyannotationreach{Matt.20.10}{Matt}{true}\chronologyannotationreach{Matt.20.11}{Matt}{true}\chronologyannotationreach{Matt.20.12}{Matt}{true}\chronologyannotationreach{Matt.20.13}{Matt}{true}\chronologyannotationreach{Matt.20.14}{Matt}{true}\chronologyannotationreach{Matt.20.15}{Matt}{true}\chronologyannotationreach{Matt.20.16}{Matt}{true}\chronologyannotationreach{Matt.20.2}{Matt}{true}\chronologyannotationreach{Matt.20.3}{Matt}{true}\chronologyannotationreach{Matt.20.4}{Matt}{true}\chronologyannotationreach{Matt.20.5}{Matt}{true}\chronologyannotationreach{Matt.20.6}{Matt}{true}\chronologyannotationreach{Matt.20.7}{Matt}{true}\chronologyannotationreach{Matt.20.8}{Matt}{true}\chronologyannotationreach{Matt.20.9}{Matt}{true}\chronologyannotationclaim{composition.gospel-of-matthew}{composition}{catholic-traditional-v1}{disputed}{about the years 64-67}{c. A.D. 64--67}; \chronologyannotationreach{Matt.20.1}{Matt}{true}\chronologyannotationreach{Matt.20.10}{Matt}{true}\chronologyannotationreach{Matt.20.11}{Matt}{true}\chronologyannotationreach{Matt.20.12}{Matt}{true}\chronologyannotationreach{Matt.20.13}{Matt}{true}\chronologyannotationreach{Matt.20.14}{Matt}{true}\chronologyannotationreach{Matt.20.15}{Matt}{true}\chronologyannotationreach{Matt.20.16}{Matt}{true}\chronologyannotationreach{Matt.20.2}{Matt}{true}\chronologyannotationreach{Matt.20.3}{Matt}{true}\chronologyannotationreach{Matt.20.4}{Matt}{true}\chronologyannotationreach{Matt.20.5}{Matt}{true}\chronologyannotationreach{Matt.20.6}{Matt}{true}\chronologyannotationreach{Matt.20.7}{Matt}{true}\chronologyannotationreach{Matt.20.8}{Matt}{true}\chronologyannotationreach{Matt.20.9}{Matt}{true}\chronologyannotationclaim{composition.gospel-of-matthew}{composition}{catholic-traditional-v1}{disputed}{about the year 50}{c. A.D. 50}.}%
  \space \chronologyannotationcomparisongroup{catholic-critical-v1}{composition}{composition-only}{\textbf{Composition (Catholic critical chronology, for comparison)}: \chronologyannotationreach{Matt.20.1}{Matt}{true}\chronologyannotationreach{Matt.20.10}{Matt}{true}\chronologyannotationreach{Matt.20.11}{Matt}{true}\chronologyannotationreach{Matt.20.12}{Matt}{true}\chronologyannotationreach{Matt.20.13}{Matt}{true}\chronologyannotationreach{Matt.20.14}{Matt}{true}\chronologyannotationreach{Matt.20.15}{Matt}{true}\chronologyannotationreach{Matt.20.16}{Matt}{true}\chronologyannotationreach{Matt.20.2}{Matt}{true}\chronologyannotationreach{Matt.20.3}{Matt}{true}\chronologyannotationreach{Matt.20.4}{Matt}{true}\chronologyannotationreach{Matt.20.5}{Matt}{true}\chronologyannotationreach{Matt.20.6}{Matt}{true}\chronologyannotationreach{Matt.20.7}{Matt}{true}\chronologyannotationreach{Matt.20.8}{Matt}{true}\chronologyannotationreach{Matt.20.9}{Matt}{true}\chronologyannotationcomparisonclaim{critical.gospel-of-matthew}{composition}{catholic-critical-v1}{catholic-critical-v1}{preferred}{passage.united-states-conference-of-catholic-bishops.new-american-bible-revised-edition.english-usccb-web-2026-09-21.matthew-introduction}{post-A.D. 70 date}{The Greek Gospel of Matthew in the NABRE introduction, post-A.D. 70 date}.}%
}
% communion-ps119: Communion antiphon: Psalm option
\expandafter\def\csname triptychchronologyannotation@communion-ps119\endcsname{%
  \chronologyannotationgroup{traditional-attribution}{disputed}{\textbf{Traditional attribution} -- disputed: \chronologyannotationreach{Ps.118.4}{Ps.118}{true}\chronologyannotationreach{Ps.118.5}{Ps.118}{true}\chronologyannotationclaim{israel.monarchy.david-traditional-era}{traditional-attribution}{catholic-traditional-v1}{disputed}{reigned from 1055 to 1015 B.C.}{David (reign in the usual chronology), B.C. 1055--1015}.}%
  \space \chronologyannotationgroup{historical-setting}{research-pending}{\textbf{Historical setting} -- Occasion date unresolved.}%
  \space \chronologyannotationgroup{composition}{preferred}{\textbf{Composition}: \chronologyannotationreach{Ps.118.4}{Ps}{true}\chronologyannotationreach{Ps.118.5}{Ps}{true}\chronologyannotationclaim{critical.psalms.latest-composition-boundary}{composition}{catholic-critical-v1}{preferred}{before the Maccabean period, around 165 B.C.; no individual psalm can be dated securely}{Before c. 165 B.C.}}%
}
% communion-jn10: Communion antiphon: John option
\expandafter\def\csname triptychchronologyannotation@communion-jn10\endcsname{%
  \chronologyannotationgroup{composition}{disputed}{\textbf{Composition} -- disputed: \chronologyannotationreach{John.10.14}{John}{true}\chronologyannotationclaim{composition.gospel-of-john}{composition}{catholic-traditional-v1}{disputed}{from the year 90 to 100 (approximately)}{c. A.D. 90--100}; \chronologyannotationreach{John.10.14}{John}{true}\chronologyannotationclaim{composition.gospel-of-john}{composition}{catholic-traditional-v1}{disputed}{96 or one of the succeeding years}{A.D. 96--100}.}%
}
% presentation: orientation-v3
\newcommand{\chronologyorientationrow}[2]{%
  \ifcsname triptychchronologyorientation@#1\endcsname
    \csname triptychchronologyorientation@#1\endcsname{#2}%
  \else
    \PackageError{triptych}{No chronology orientation row for #1}{}%
  \fi}
\expandafter\def\csname triptychchronologyorientation@row-dca539dc5ada8d26\endcsname#1{%
Entrance antiphon: Ps. 37:39--40 (Heb.; Vg. 36) & Historical setting & Occasion date unresolved\\*[.1em]%
}
\expandafter\def\csname triptychchronologyorientation@row-0fd3ac149bd11fa5\endcsname#1{%
 & Composition: Psalms & \chronologyannotationreach{Ps.36.39}{Ps}{true}\chronologyannotationreach{Ps.36.40}{Ps}{true}\chronologyannotationclaim{critical.psalms.latest-composition-boundary}{composition}{catholic-critical-v1}{preferred}{before the Maccabean period, around 165 B.C.; no individual psalm can be dated securely}{Before c.~165 B.C.}\\[.3em]%
}
\expandafter\def\csname triptychchronologyorientation@row-40514a2e4f93a2b4\endcsname#1{%
Communion antiphon: Psalm option: Ps. 119:4--5 (Heb.; Vg. 118) & Traditional attribution: David's reign in the usual chronology; Hebron, then Jerusalem & Disputed: \chronologyannotationreach{Ps.118.4}{Ps.118}{true}\chronologyannotationreach{Ps.118.5}{Ps.118}{true}\chronologyannotationclaim{israel.monarchy.david-traditional-era}{traditional-attribution}{catholic-traditional-v1}{disputed}{reigned from 1055 to 1015 B.C.}{B.C.~1055--1015 (Corbett also reports writers dating the reign 30--50 years later)}\\*[.1em]%
}
\expandafter\def\csname triptychchronologyorientation@row-ef1c0be7f588413e\endcsname#1{%
 & Historical setting & Occasion date unresolved\\*[.1em]%
}
\expandafter\def\csname triptychchronologyorientation@row-489e7316a0e583d7\endcsname#1{%
 & Composition: Psalms & \chronologyannotationreach{Ps.118.4}{Ps}{true}\chronologyannotationreach{Ps.118.5}{Ps}{true}\chronologyannotationclaim{critical.psalms.latest-composition-boundary}{composition}{catholic-critical-v1}{preferred}{before the Maccabean period, around 165 B.C.; no individual psalm can be dated securely}{Before c.~165 B.C.}\\[.3em]%
}
\expandafter\def\csname triptychchronologyorientation@row-2764665a12e4f1f9\endcsname#1{%
Responsorial psalm: Ps. 145:2--3, 8--9, 17--18 (Heb.; Vg. 144), Ps. 145:18a (Heb.; Vg. 144) & Traditional attribution: David's reign in the usual chronology; Hebron, then Jerusalem & Disputed: \chronologyannotationreach{Ps.144.17}{Ps.107 Ps.144}{true}\chronologyannotationreach{Ps.144.18}{Ps.107 Ps.144}{true}\chronologyannotationreach{Ps.144.2}{Ps.107 Ps.144}{true}\chronologyannotationreach{Ps.144.3}{Ps.107 Ps.144}{true}\chronologyannotationreach{Ps.144.8}{Ps.107 Ps.144}{true}\chronologyannotationreach{Ps.144.9}{Ps.107 Ps.144}{true}\chronologyannotationclaim{israel.monarchy.david-traditional-era}{traditional-attribution}{catholic-traditional-v1}{disputed}{reigned from 1055 to 1015 B.C.}{B.C.~1055--1015 (Corbett also reports writers dating the reign 30--50 years later)}\\*[.1em]%
}
\expandafter\def\csname triptychchronologyorientation@row-6e38bd9e0e57375d\endcsname#1{%
 & Historical setting & Occasion date unresolved\\*[.1em]%
}
\expandafter\def\csname triptychchronologyorientation@row-632418cf547f5d49\endcsname#1{%
 & Composition: Psalms & \chronologyannotationreach{Ps.144.17}{Ps}{true}\chronologyannotationreach{Ps.144.18}{Ps}{true}\chronologyannotationreach{Ps.144.2}{Ps}{true}\chronologyannotationreach{Ps.144.3}{Ps}{true}\chronologyannotationreach{Ps.144.8}{Ps}{true}\chronologyannotationreach{Ps.144.9}{Ps}{true}\chronologyannotationclaim{critical.psalms.latest-composition-boundary}{composition}{catholic-critical-v1}{preferred}{before the Maccabean period, around 165 B.C.; no individual psalm can be dated securely}{Before c.~165 B.C.}\\[.3em]%
}
\expandafter\def\csname triptychchronologyorientation@row-e2ada12f5ff3422d\endcsname#1{%
First reading: Is. 55:6--9 & Traditional attribution: Isaias (ministry in Souvay's traditional account); Judah and Jerusalem & \chronologyannotationreach{Is.55.6}{Is}{true}\chronologyannotationreach{Is.55.7}{Is}{true}\chronologyannotationreach{Is.55.8}{Is}{true}\chronologyannotationreach{Is.55.9}{Is}{true}\chronologyannotationclaim{israel.prophets.isaias-traditional-era}{traditional-attribution}{catholic-traditional-v1}{preferred}{from the closing year of Ozias, King of Juda (740 B.C.); no prophecies appear to be later than 701 B.C.}{B.C.~740--701}\\*[.1em]%
}
\expandafter\def\csname triptychchronologyorientation@row-1d46ab7b264ff05e\endcsname#1{%
 & Prophecy given: The prophetic activity behind Isaiah 40-55 in the NABRE introduction; exiles from Babylon; restoration to Jerusalem & \chronologyannotationreach{Is.55.6}{Is.40 Is.41 Is.42 Is.43 Is.44 Is.45 Is.46 Is.47 Is.48 Is.49 Is.50 Is.51 Is.52 Is.53 Is.54 Is.55}{true}\chronologyannotationreach{Is.55.7}{Is.40 Is.41 Is.42 Is.43 Is.44 Is.45 Is.46 Is.47 Is.48 Is.49 Is.50 Is.51 Is.52 Is.53 Is.54 Is.55}{true}\chronologyannotationreach{Is.55.8}{Is.40 Is.41 Is.42 Is.43 Is.44 Is.45 Is.46 Is.47 Is.48 Is.49 Is.50 Is.51 Is.52 Is.53 Is.54 Is.55}{true}\chronologyannotationreach{Is.55.9}{Is.40 Is.41 Is.42 Is.43 Is.44 Is.45 Is.46 Is.47 Is.48 Is.49 Is.50 Is.51 Is.52 Is.53 Is.54 Is.55}{true}\chronologyannotationclaim{israel.exile.second-isaias-oracles}{prophecy-given}{catholic-critical-v1}{preferred}{toward the end of the Babylonian exile}{Toward the end of the Babylonian exile}\\[.3em]%
}
\expandafter\def\csname triptychchronologyorientation@row-882de3ad7e508ff5\endcsname#1{%
Gospel: Matt. 20:1--16a & Narrated event & Date unresolved\\*[.1em]%
}
\expandafter\def\csname triptychchronologyorientation@row-c947b22e36134504\endcsname#1{%
 & Composition: Matthew; Semitic original associated with Palestine & Disputed: \chronologyannotationreach{Matt.20.1}{Matt}{true}\chronologyannotationreach{Matt.20.10}{Matt}{true}\chronologyannotationreach{Matt.20.11}{Matt}{true}\chronologyannotationreach{Matt.20.12}{Matt}{true}\chronologyannotationreach{Matt.20.13}{Matt}{true}\chronologyannotationreach{Matt.20.14}{Matt}{true}\chronologyannotationreach{Matt.20.15}{Matt}{true}\chronologyannotationreach{Matt.20.16}{Matt}{true}\chronologyannotationreach{Matt.20.2}{Matt}{true}\chronologyannotationreach{Matt.20.3}{Matt}{true}\chronologyannotationreach{Matt.20.4}{Matt}{true}\chronologyannotationreach{Matt.20.5}{Matt}{true}\chronologyannotationreach{Matt.20.6}{Matt}{true}\chronologyannotationreach{Matt.20.7}{Matt}{true}\chronologyannotationreach{Matt.20.8}{Matt}{true}\chronologyannotationreach{Matt.20.9}{Matt}{true}\chronologyannotationclaim{composition.gospel-of-matthew}{composition}{catholic-traditional-v1}{disputed}{about A.D. 38-45}{c.~A.D.~38--45}; \chronologyannotationreach{Matt.20.1}{Matt}{true}\chronologyannotationreach{Matt.20.10}{Matt}{true}\chronologyannotationreach{Matt.20.11}{Matt}{true}\chronologyannotationreach{Matt.20.12}{Matt}{true}\chronologyannotationreach{Matt.20.13}{Matt}{true}\chronologyannotationreach{Matt.20.14}{Matt}{true}\chronologyannotationreach{Matt.20.15}{Matt}{true}\chronologyannotationreach{Matt.20.16}{Matt}{true}\chronologyannotationreach{Matt.20.2}{Matt}{true}\chronologyannotationreach{Matt.20.3}{Matt}{true}\chronologyannotationreach{Matt.20.4}{Matt}{true}\chronologyannotationreach{Matt.20.5}{Matt}{true}\chronologyannotationreach{Matt.20.6}{Matt}{true}\chronologyannotationreach{Matt.20.7}{Matt}{true}\chronologyannotationreach{Matt.20.8}{Matt}{true}\chronologyannotationreach{Matt.20.9}{Matt}{true}\chronologyannotationclaim{composition.gospel-of-matthew}{composition}{catholic-traditional-v1}{disputed}{about the year 40-42}{c.~A.D.~40--42}; \chronologyannotationreach{Matt.20.1}{Matt}{true}\chronologyannotationreach{Matt.20.10}{Matt}{true}\chronologyannotationreach{Matt.20.11}{Matt}{true}\chronologyannotationreach{Matt.20.12}{Matt}{true}\chronologyannotationreach{Matt.20.13}{Matt}{true}\chronologyannotationreach{Matt.20.14}{Matt}{true}\chronologyannotationreach{Matt.20.15}{Matt}{true}\chronologyannotationreach{Matt.20.16}{Matt}{true}\chronologyannotationreach{Matt.20.2}{Matt}{true}\chronologyannotationreach{Matt.20.3}{Matt}{true}\chronologyannotationreach{Matt.20.4}{Matt}{true}\chronologyannotationreach{Matt.20.5}{Matt}{true}\chronologyannotationreach{Matt.20.6}{Matt}{true}\chronologyannotationreach{Matt.20.7}{Matt}{true}\chronologyannotationreach{Matt.20.8}{Matt}{true}\chronologyannotationreach{Matt.20.9}{Matt}{true}\chronologyannotationclaim{composition.gospel-of-matthew}{composition}{catholic-traditional-v1}{disputed}{the years 40-45}{A.D.~40--45}; \chronologyannotationreach{Matt.20.1}{Matt}{true}\chronologyannotationreach{Matt.20.10}{Matt}{true}\chronologyannotationreach{Matt.20.11}{Matt}{true}\chronologyannotationreach{Matt.20.12}{Matt}{true}\chronologyannotationreach{Matt.20.13}{Matt}{true}\chronologyannotationreach{Matt.20.14}{Matt}{true}\chronologyannotationreach{Matt.20.15}{Matt}{true}\chronologyannotationreach{Matt.20.16}{Matt}{true}\chronologyannotationreach{Matt.20.2}{Matt}{true}\chronologyannotationreach{Matt.20.3}{Matt}{true}\chronologyannotationreach{Matt.20.4}{Matt}{true}\chronologyannotationreach{Matt.20.5}{Matt}{true}\chronologyannotationreach{Matt.20.6}{Matt}{true}\chronologyannotationreach{Matt.20.7}{Matt}{true}\chronologyannotationreach{Matt.20.8}{Matt}{true}\chronologyannotationreach{Matt.20.9}{Matt}{true}\chronologyannotationclaim{composition.gospel-of-matthew}{composition}{catholic-traditional-v1}{disputed}{about the year 60-68}{c.~A.D.~60--68}; \chronologyannotationreach{Matt.20.1}{Matt}{true}\chronologyannotationreach{Matt.20.10}{Matt}{true}\chronologyannotationreach{Matt.20.11}{Matt}{true}\chronologyannotationreach{Matt.20.12}{Matt}{true}\chronologyannotationreach{Matt.20.13}{Matt}{true}\chronologyannotationreach{Matt.20.14}{Matt}{true}\chronologyannotationreach{Matt.20.15}{Matt}{true}\chronologyannotationreach{Matt.20.16}{Matt}{true}\chronologyannotationreach{Matt.20.2}{Matt}{true}\chronologyannotationreach{Matt.20.3}{Matt}{true}\chronologyannotationreach{Matt.20.4}{Matt}{true}\chronologyannotationreach{Matt.20.5}{Matt}{true}\chronologyannotationreach{Matt.20.6}{Matt}{true}\chronologyannotationreach{Matt.20.7}{Matt}{true}\chronologyannotationreach{Matt.20.8}{Matt}{true}\chronologyannotationreach{Matt.20.9}{Matt}{true}\chronologyannotationclaim{composition.gospel-of-matthew}{composition}{catholic-traditional-v1}{disputed}{about the years 64-67}{c.~A.D.~64--67}; \chronologyannotationreach{Matt.20.1}{Matt}{true}\chronologyannotationreach{Matt.20.10}{Matt}{true}\chronologyannotationreach{Matt.20.11}{Matt}{true}\chronologyannotationreach{Matt.20.12}{Matt}{true}\chronologyannotationreach{Matt.20.13}{Matt}{true}\chronologyannotationreach{Matt.20.14}{Matt}{true}\chronologyannotationreach{Matt.20.15}{Matt}{true}\chronologyannotationreach{Matt.20.16}{Matt}{true}\chronologyannotationreach{Matt.20.2}{Matt}{true}\chronologyannotationreach{Matt.20.3}{Matt}{true}\chronologyannotationreach{Matt.20.4}{Matt}{true}\chronologyannotationreach{Matt.20.5}{Matt}{true}\chronologyannotationreach{Matt.20.6}{Matt}{true}\chronologyannotationreach{Matt.20.7}{Matt}{true}\chronologyannotationreach{Matt.20.8}{Matt}{true}\chronologyannotationreach{Matt.20.9}{Matt}{true}\chronologyannotationclaim{composition.gospel-of-matthew}{composition}{catholic-traditional-v1}{disputed}{about the year 50}{c.~A.D.~50}\\*[.1em]%
}
\expandafter\def\csname triptychchronologyorientation@row-3d55a60db0e5ee05\endcsname#1{%
 & Composition, Catholic critical comparison: Greek Matthew; Antioch suggested & \chronologyannotationreach{Matt.20.1}{Matt}{true}\chronologyannotationreach{Matt.20.10}{Matt}{true}\chronologyannotationreach{Matt.20.11}{Matt}{true}\chronologyannotationreach{Matt.20.12}{Matt}{true}\chronologyannotationreach{Matt.20.13}{Matt}{true}\chronologyannotationreach{Matt.20.14}{Matt}{true}\chronologyannotationreach{Matt.20.15}{Matt}{true}\chronologyannotationreach{Matt.20.16}{Matt}{true}\chronologyannotationreach{Matt.20.2}{Matt}{true}\chronologyannotationreach{Matt.20.3}{Matt}{true}\chronologyannotationreach{Matt.20.4}{Matt}{true}\chronologyannotationreach{Matt.20.5}{Matt}{true}\chronologyannotationreach{Matt.20.6}{Matt}{true}\chronologyannotationreach{Matt.20.7}{Matt}{true}\chronologyannotationreach{Matt.20.8}{Matt}{true}\chronologyannotationreach{Matt.20.9}{Matt}{true}\chronologyannotationcomparisonclaim{critical.gospel-of-matthew}{composition}{catholic-critical-v1}{catholic-critical-v1}{preferred}{passage.united-states-conference-of-catholic-bishops.new-american-bible-revised-edition.english-usccb-web-2026-09-21.matthew-introduction}{post-A.D. 70 date}{Post-A.D.~70 date}\\[.3em]%
}
\expandafter\def\csname triptychchronologyorientation@row-fc88dd0b7b1b3ba4\endcsname#1{%
Communion antiphon: John option: John 10:14 & Composition: John; Ephesus traditionally & Disputed: \chronologyannotationreach{John.10.14}{John}{true}\chronologyannotationclaim{composition.gospel-of-john}{composition}{catholic-traditional-v1}{disputed}{from the year 90 to 100 (approximately)}{c.~A.D.~90--100}; \chronologyannotationreach{John.10.14}{John}{true}\chronologyannotationclaim{composition.gospel-of-john}{composition}{catholic-traditional-v1}{disputed}{96 or one of the succeeding years}{A.D.~96--100}\\[.3em]%
}
\expandafter\def\csname triptychchronologyorientation@row-b7c0e93d90a711f8\endcsname#1{%
Gospel acclamation: Cf. Acts 16:14b & Narrated event: St Paul's Second Missionary Journey & Disputed: \chronologyannotationreach{Acts.16.14}{Acts.15.36- Acts.16 Acts.17 Acts.18.1-22}{true}\chronologyannotationclaim{apostolic-age.saint-paul-second-mission}{narrated-event}{catholic-traditional-v1}{disputed}{second mission, 50-53}{A.D.~50--53}; \chronologyannotationreach{Acts.16.14}{Acts.15.36- Acts.16 Acts.17 Acts.18.1-22}{true}\chronologyannotationclaim{apostolic-age.saint-paul-second-mission}{narrated-event}{catholic-traditional-v1}{disputed}{began about three years previously, in 51 ... probably in 54}{A.D.~51--54}\\*[.1em]%
}
\expandafter\def\csname triptychchronologyorientation@row-8945f74134c1b01a\endcsname#1{%
 & Composition: The Acts of the Apostles & \chronologyannotationreach{Acts.16.14}{Acts}{true}\chronologyannotationclaim{composition.acts-of-the-apostles}{composition}{catholic-traditional-v1}{preferred}{the year 64 is the most probable date}{c.~A.D.~64}\\[.3em]%
}
\expandafter\def\csname triptychchronologyorientation@row-4f1fbee7c004e5bd\endcsname#1{%
Second reading: Phil. 1:20c--24, 27a & Composition: Philippians; Rome to Philippi in the cited account & Disputed: \chronologyannotationreach{Phil.1.20}{Phil}{true}\chronologyannotationreach{Phil.1.21}{Phil}{true}\chronologyannotationreach{Phil.1.22}{Phil}{true}\chronologyannotationreach{Phil.1.23}{Phil}{true}\chronologyannotationreach{Phil.1.24}{Phil}{true}\chronologyannotationreach{Phil.1.27}{Phil}{true}\chronologyannotationclaim{composition.epistle-to-the-philippians}{composition}{catholic-traditional-v1}{disputed}{(Philemon; Colossians; Ephesians; Philippians), 61}{A.D.~61}; \chronologyannotationreach{Phil.1.20}{Phil}{true}\chronologyannotationreach{Phil.1.21}{Phil}{true}\chronologyannotationreach{Phil.1.22}{Phil}{true}\chronologyannotationreach{Phil.1.23}{Phil}{true}\chronologyannotationreach{Phil.1.24}{Phil}{true}\chronologyannotationreach{Phil.1.27}{Phil}{true}\chronologyannotationclaim{composition.epistle-to-the-philippians}{composition}{catholic-traditional-v1}{disputed}{at Rome (A.D. 62-64)}{A.D.~62--64}\\[.3em]%
}

% Generated. Do not edit.
% Complete orientation-v3 table; wording, dates and order belong to the corpus.
\begin{chronologysharedtable}
\chronologyorientationrow{row-dca539dc5ada8d26}{}
\chronologyorientationrow{row-0fd3ac149bd11fa5}{}
\midrule
\chronologyorientationrow{row-17feb57cbca8d8c6}{}
\chronologyorientationrow{row-68e7d382f94a3249}{}
\chronologyorientationrow{row-98b280443fe96ee4}{}
\midrule
\chronologyorientationrow{row-5d22841015565a14}{}
\chronologyorientationrow{row-76516bdba892ff47}{}
\chronologyorientationrow{row-bbf9c28308627449}{}
\midrule
\chronologyorientationrow{row-e2ada12f5ff3422d}{}
\chronologyorientationrow{row-1d46ab7b264ff05e}{}
\midrule
\chronologyorientationrow{row-882de3ad7e508ff5}{}
\chronologyorientationrow{row-c947b22e36134504}{}
\chronologyorientationrow{row-3d55a60db0e5ee05}{}
\midrule
\chronologyorientationrow{row-fc88dd0b7b1b3ba4}{}
\midrule
\chronologyorientationrow{row-b7c0e93d90a711f8}{}
\chronologyorientationrow{row-8945f74134c1b01a}{}
\midrule
\chronologyorientationrow{row-4f1fbee7c004e5bd}{}
\end{chronologysharedtable}
{\footnotesize Psalm numbering: Vg. = Vulgate; Heb. = Hebrew.\par}


\subsection{Evidence for each date}
\label{sec:date-location-evidence}
{\small For each passage, as the Scriptural Date and Location table names it: what is dated, how the corpus weighs each answer, the source's own words, and the sources cited. A subject the corpus does not date is reported as such and given no date.\par}
{\small
\par\smallskip\noindent\textit{Entrance antiphon: Ps. 37:39--40 (Heb.; Vg. 36)}\par\nobreak
\begin{itemize}[leftmargin=1.2em,itemsep=1pt,topsep=1pt,parsep=0pt]
\item \textbf{Historical setting}. No date: the chronology corpus has not yet dated this (research pending).
\item \textbf{Composition} (preferred). Psalms: Before c.~165 B.C. (``before the Maccabean period, around 165 B.C.; no individual psalm can be dated securely''; Catholic critical chronology; New American Bible, Revised Edition --- USCCB web state (2026), Psalms introduction).
\end{itemize}
\par\smallskip\noindent\textit{Communion antiphon: Psalm option: Ps. 119:4--5 (Heb.; Vg. 118)}\par\nobreak
\begin{itemize}[leftmargin=1.2em,itemsep=1pt,topsep=1pt,parsep=0pt]
\item \textbf{Traditional attribution} (disputed). David's reign: B.C.~1055--1015 (``reigned from 1055 to 1015 B.C.''; The Catholic Encyclopedia, Volume 4 (New York, 1908), article 'David, King,' printed p. 642, opening paragraph). Place: Hebron, then Jerusalem (The Holy Bible, Douay-Rheims (Challoner revision), Project Gutenberg eBook 1581: 2 Kings 5:5).
\item \textbf{Historical setting}. As under Entrance antiphon: Ps. 37:39--40 (Heb.; Vg. 36).
\item \textbf{Composition}. As under Entrance antiphon: Ps. 37:39--40 (Heb.; Vg. 36).
\end{itemize}
\par\smallskip\noindent\textit{Responsorial psalm: Ps. 145:2--3, 8--9, 17--18 (Heb.; Vg. 144), Ps. 145:18a (Heb.; Vg. 144)}\par\nobreak
\begin{itemize}[leftmargin=1.2em,itemsep=1pt,topsep=1pt,parsep=0pt]
\item \textbf{Traditional attribution}. As under Communion antiphon: Psalm option: Ps. 119:4--5 (Heb.; Vg. 118).
\item \textbf{Historical setting}. As under Entrance antiphon: Ps. 37:39--40 (Heb.; Vg. 36).
\item \textbf{Composition}. As under Entrance antiphon: Ps. 37:39--40 (Heb.; Vg. 36).
\end{itemize}
\par\smallskip\noindent\textit{First reading: Is. 55:6--9}\par\nobreak
\begin{itemize}[leftmargin=1.2em,itemsep=1pt,topsep=1pt,parsep=0pt]
\item \textbf{Traditional attribution} (preferred). Isaias (ministry in Souvay's traditional account): B.C.~740--701 (``from the closing year of Ozias, King of Juda (740 B.C.); no prophecies appear to be later than 701 B.C.''; The Catholic Encyclopedia, Volume 8 (New York, 1910)). Place: Judah and Jerusalem (The Holy Bible, Douay-Rheims (Challoner revision), Project Gutenberg eBook 1581: Is. 1:1).
\item \textbf{Prophecy given} (preferred). The prophetic activity behind Isaiah 40-55 in the NABRE introduction: Toward the end of the Babylonian exile (Catholic critical chronology; New American Bible, Revised Edition: USCCB web state (2026), Isaiah introduction, especially the final two prose paragraphs before the outline). Place: exiles from Babylon; restoration to Jerusalem (The Holy Bible, Douay-Rheims (Challoner revision), Project Gutenberg eBook 1581: Is. 48:20, 52:9).
\end{itemize}
\par\smallskip\noindent\textit{Gospel: Matt. 20:1--16a}\par\nobreak
\begin{itemize}[leftmargin=1.2em,itemsep=1pt,topsep=1pt,parsep=0pt]
\item \textbf{Narrated event}. No date: the chronology corpus has not yet dated this (research pending).
\item \textbf{Composition} (disputed). The Gospel of St Matthew: c.~A.D.~38--45 (``about A.D. 38-45''; The Catholic Encyclopedia, Volume 10 (1911); The Catholic Encyclopedia, Volume 14 (New York, 1912)); c.~A.D.~40--42 (``about the year 40-42''; The Catholic Encyclopedia, Volume 10 (1911)); A.D.~40--45 (``the years 40-45''; The Catholic Encyclopedia, Volume 10 (1911)); c.~A.D.~60--68 (``about the year 60-68''; The Catholic Encyclopedia, Volume 10 (1911)); c.~A.D.~64--67 (``about the years 64-67''; The Catholic Encyclopedia, Volume 10 (1911)); c.~A.D.~50 (``about the year 50''; The Catholic Encyclopedia, Volume 14 (New York, 1912)). Place: Semitic original associated with Palestine (The Catholic Encyclopedia, Volume 10 (1911)).
\item \textbf{Composition, Catholic critical comparison} (preferred). The Greek Gospel of Matthew in the NABRE introduction: Post-A.D.~70 date (Catholic critical chronology; New American Bible, Revised Edition: USCCB web state (2026), Matthew introduction, final five prose paragraphs on authorship, sources, composition and place; narrative overview of chapters 19-25). Place: Antioch suggested (New American Bible, Revised Edition: USCCB web state (2026), Matthew introduction, final five prose paragraphs on authorship, sources, composition and place; narrative overview of chapters 19-25).
\end{itemize}
\par\smallskip\noindent\textit{Communion antiphon: John option: John 10:14}\par\nobreak
\begin{itemize}[leftmargin=1.2em,itemsep=1pt,topsep=1pt,parsep=0pt]
\item \textbf{Composition} (disputed). The Gospel of St John: c.~A.D.~90--100 (``from the year 90 to 100 (approximately)''; The Catholic Encyclopedia, Volume 14 (New York, 1912)); A.D.~96--100 (``96 or one of the succeeding years''; The Catholic Encyclopedia, Volume 8 (New York, 1910)). Place: Ephesus traditionally (The Catholic Encyclopedia, Volume 8 (New York, 1910)).
\end{itemize}
\par\smallskip\noindent\textit{Gospel acclamation: Cf. Acts 16:14b}\par\nobreak
\begin{itemize}[leftmargin=1.2em,itemsep=1pt,topsep=1pt,parsep=0pt]
\item \textbf{Narrated event} (disputed). St Paul's Second Missionary Journey: A.D.~50--53 (``second mission, 50-53''; The Catholic Encyclopedia, Volume 11 (New York, 1911)); A.D.~51--54 (``began about three years previously, in 51 ... probably in 54''; The Catholic Encyclopedia, Volume 3 (New York, 1908)).
\item \textbf{Composition} (preferred). The Acts of the Apostles: c.~A.D.~64 (``the year 64 is the most probable date''; The Catholic Encyclopedia, Volume 1 (New York, 1907)).
\end{itemize}
\par\smallskip\noindent\textit{Second reading: Phil. 1:20c--24, 27a}\par\nobreak
\begin{itemize}[leftmargin=1.2em,itemsep=1pt,topsep=1pt,parsep=0pt]
\item \textbf{Composition} (disputed). The Epistle to the Philippians: A.D.~61 (``(Philemon; Colossians; Ephesians; Philippians), 61''; The Catholic Encyclopedia, Volume 11 (New York, 1911)); A.D.~62--64 (``at Rome (A.D. 62-64)''; The Catholic Encyclopedia, Volume 12 (1911)). Place: Rome to Philippi in the cited account (The Catholic Encyclopedia, Volume 12 (1911); The Holy Bible, Douay-Rheims (Challoner revision), Project Gutenberg eBook 1581: Phil. 1:1).
\end{itemize}
\par}
